\section{4. Especificación de requisitos de software
(SRS)}\label{especificaciuxf3n-de-requisitos-de-software-srs}

Estructura adaptada de ISO/IEC/IEEE 29148:2018. Cada requisito tiene un
identificador único que se utiliza para la trazabilidad.

\subsection{4.1 Propósito, alcance y audiencia del
SRS}\label{propuxf3sito-alcance-y-audiencia-del-srs}

\textbf{Propósito:} este SRS describe qué debe hacer el Sistema de
Órdenes de Servicio y Citas y con qué calidad, para que el cliente lo
valide y lo firme en la Review 2, y para que el equipo lo use como base
del diseño, la implementación y las pruebas.

\textbf{Alcance:} cubre el ciclo de atención del taller (citas,
recepción, diagnóstico, cotización y autorización, refacciones, cobro y
entrega) y el portal web del cliente, según el alcance de la sección
3.3.

\textbf{Audiencia:} los dueños y el gerente del taller (aprobación), el
personal que usará el sistema (validación de sus flujos), el equipo de
desarrollo y el docente evaluador.

\textbf{Convenciones:} cada requisito tiene un identificador único: RF
(funcional), RNF (no funcional), RN (regla de negocio), HU (historia de
usuario) y CU (caso de uso). Estos identificadores se usan en la matriz
de trazabilidad (5.3).

\subsection{4.2 Glosario del dominio}\label{glosario-del-dominio}

{\def\LTcaptype{none} % do not increment counter
\begin{longtable}[]{@{}
  >{\raggedright\arraybackslash}p{(\linewidth - 4\tabcolsep) * \real{0.2349}}
  >{\raggedright\arraybackslash}p{(\linewidth - 4\tabcolsep) * \real{0.5342}}
  >{\raggedright\arraybackslash}p{(\linewidth - 4\tabcolsep) * \real{0.2309}}@{}}
\toprule\noalign{}
\begin{minipage}[b]{\linewidth}\raggedright
\textbf{Término}
\end{minipage} & \begin{minipage}[b]{\linewidth}\raggedright
\textbf{Definición}
\end{minipage} & \begin{minipage}[b]{\linewidth}\raggedright
\textbf{Sinónimos que usa el cliente}
\end{minipage} \\
\midrule\noalign{}
\endhead
\bottomrule\noalign{}
\endlastfoot
Autorización & Aceptación de una cotización por el cliente; registra
fecha, hora, IP, monto y el código de verificación. & Aprobación, firma
digital \\
Bahía & Espacio de trabajo donde se repara un vehículo. El taller tiene
4 y cada una aloja un vehículo a la vez. & --- \\
Checklist de inspección & Lista de verificación del estado de
carrocería, interiores y accesorios al recibir el vehículo. & Lista de
inspección \\
Cita & Solicitud para atender un vehículo en una fecha y hora; queda
«Pendiente» hasta que la recepcionista la confirma. & --- \\
Código 2FA & Código de 6 dígitos, de un solo uso y vigente 5 minutos,
que se envía al correo del cliente para confirmar una autorización. &
Código de verificación, doble verificación \\
Compra & Adquisición de refacciones a un proveedor, respaldada por un
ticket. & Ticket de compra \\
Concepto & Renglón de una cotización: mano de obra, refacción o insumo.
Puede ser opcional. & Rubro \\
Cotización & Costo desglosado de mano de obra, refacciones e insumos de
una orden; vigente 5 días naturales. & Presupuesto \\
Diagnóstico & Descripción técnica de la falla y de las reparaciones
necesarias que registra el jefe de mecánicos. & --- \\
Enlace seguro & Liga web única que recibe el cliente para consultar su
vehículo y autorizar cotizaciones sin usuario ni contraseña. & Token de
acceso \\
Entrega a crédito & Entrega del vehículo con saldo pendiente; solo para
flotillas registradas y con autorización del gerente. & Crédito \\
Flotilla & Empresa cliente con varios vehículos que puede recibir
entrega a crédito. & Flotilla empresarial \\
Folio & Identificador único de una orden de servicio. & --- \\
Fotografía de recepción & Evidencia digital del estado del vehículo al
llegar. Son obligatorias cuatro: frente, trasera, costado izquierdo y
costado derecho. & --- \\
Garantía & Cobertura de 30 días sobre la mano de obra, contados desde la
entrega del vehículo. & --- \\
Insumo & Material consumible que se usa en la reparación, como aceite o
líquidos. & Consumible \\
Jefe de mecánicos & Líder técnico: diagnostica, cotiza, asigna bahías y
mecánicos y solicita refacciones desde la tableta. & Líder técnico \\
Mecánico responsable & Mecánico titular de una orden; otros mecánicos
pueden colaborar sin cambiar la titularidad. & Mecánico principal \\
Orden de servicio & Registro de todo el trabajo sobre un vehículo, desde
la recepción hasta la entrega. & Orden \\
Proveedor & Refaccionaria local a la que el almacén compra piezas. &
Refaccionaria \\
Recepción vehicular & Registro de la llegada del vehículo: kilometraje,
nivel de combustible, checklist y fotografías. & Check-in, registro de
entrada \\
Re-cotización & Nueva versión de la cotización cuando los trabajos
adicionales aumentan el monto más de 10 \%. & --- \\
Refacción & Pieza que se instala en el vehículo. & Pieza, repuesto \\
Saldo pendiente & Diferencia entre el total autorizado de la orden y lo
que el cliente ha pagado. & --- \\
Solicitud de refacción & Pedido de piezas que el jefe de mecánicos hace
al almacén para una orden. & Requisición \\
\end{longtable}
}

\subsection{4.3 Descripción general}\label{descripciuxf3n-general}

\subsubsection{4.3.1 Usuarios y sus
características}\label{usuarios-y-sus-caracteruxedsticas}

El sistema será utilizado por distintos perfiles dentro de la estructura
operativa del negocio y la parte gerencial del taller mecánico, además
de los clientes finales del negocio. La siguiente tabla detalla las
características, nivel técnico y contexto de uso:

{\def\LTcaptype{none} % do not increment counter
\begin{longtable}[]{@{}
  >{\raggedright\arraybackslash}p{(\linewidth - 6\tabcolsep) * \real{0.2138}}
  >{\raggedright\arraybackslash}p{(\linewidth - 6\tabcolsep) * \real{0.2990}}
  >{\raggedright\arraybackslash}p{(\linewidth - 6\tabcolsep) * \real{0.2138}}
  >{\raggedright\arraybackslash}p{(\linewidth - 6\tabcolsep) * \real{0.2734}}@{}}
\toprule\noalign{}
\begin{minipage}[b]{\linewidth}\raggedright
\textbf{Tipo de usuario}
\end{minipage} & \begin{minipage}[b]{\linewidth}\raggedright
\textbf{Tareas principales}
\end{minipage} & \begin{minipage}[b]{\linewidth}\raggedright
\textbf{Experiencia tecnológica}
\end{minipage} & \begin{minipage}[b]{\linewidth}\raggedright
\textbf{Dispositivo y contexto de uso}
\end{minipage} \\
\midrule\noalign{}
\endhead
\bottomrule\noalign{}
\endlastfoot
Recepcionista & Registro de clientes/vehículos, check-in con fotos,
apertura de órdenes, confirmación de citas web, cobranza y entrega. &
Media: Manejo fluido de navegación web, correo electrónico y paquetes de
oficina. & Computadora de escritorio en mostrador de recepción. Entorno
de oficina con teclado/mouse. \\
Jefe de Mecánicos (Líder Técnico) & Diagnóstico técnico, elaboración de
cotizaciones, asignación de bahías/mecánicos, requisición de refacciones
y sellado final. & Baja--Media: Familiarizado con smartphones
personales; poca experiencia con sistemas de gestión. & Tabletas
táctiles exclusivas de taller (1 o 2 dispositivos). Entorno con ruido,
polvo y manipulación con manos sucias/grasosas. \\
Mecánico en Bahía & Ejecución de reparaciones físicas, notificación de
hallazgos adicionales o solicitud de refacciones al Jefe de Mecánicos. &
Baja: Uso de teléfonos inteligentes para uso personal cotidiano. & Sin
dispositivo asignado directo; interactúa a través del Jefe de Mecánicos
para evitar deterioro de hardware \\
Encargado de Almacén & Atender solicitudes de piezas, consultar
directorio de proveedores locales, registrar compras/tickets e historial
de insumos. & Baja--Media\textbf{:} Uso básico de computadora para
consulta de inventarios. & Computadora de escritorio o terminal dedicada
en el área de almacén/refacciones. \\
Gerente (Dueño) & Monitoreo del Dashboard de KPIs (ocupación, ingresos),
aprobación de entregas a crédito a flotillas y excepciones. & Media:
Manejo de smartphones y computadoras para administración del negocio. &
Laptop personal o dispositivo móvil (smartphone/tableta) con acceso
remoto/local \\
Cliente & Solicitud de citas en línea, seguimiento en vivo del servicio
y aprobación/rechazo de cotizaciones mediante código 2FA. & Variable
(Baja a Alta): Desde usuarios sin conocimientos técnicos hasta usuarios
avanzados. & Dispositivo móvil personal (iOS/Android) vía navegador web
progresivo (PWA) sin instalar aplicaciones. \\
\end{longtable}
}

\subsubsection{4.3.2 Restricciones, supuestos y
dependencias}\label{restricciones-supuestos-y-dependencias}

\begin{enumerate}
\def\labelenumi{\arabic{enumi}.}
\item
  \textbf{Restricciones:}
\end{enumerate}

\begin{itemize}
\item
  Presupuesto y Licenciamiento: El desarrollo debe realizarse
  estrictamente con herramientas, lenguajes y frameworks de código
  abierto (\emph{Open Source} como Python y FastAPI), garantizando cero
  costos de licencias o rentas de software.
\item
  Infraestructura Física: El taller cuenta únicamente con 4 bahías de
  trabajo físicas (capacidad estricta de 1 vehículo por bahía) y un
  estacionamiento con capacidad máxima de 5 vehículos en espera.
\item
  Dispositivos Hardware: La captura digital en taller se limita a máximo
  2 tabletas táctiles centralizadas por el Jefe de Mecánicos para
  proteger la inversión y evitar daños.
\end{itemize}

\begin{itemize}
\item
  Alcance del Vehículo\textbf{:} La plataforma atenderá exclusivamente
  vehículos ligeros (automóviles particulares, motocicletas y
  camionetas), excluyendo maquinaria pesada o industrial.
\end{itemize}

\begin{enumerate}
\def\labelenumi{\arabic{enumi}.}
\setcounter{enumi}{1}
\item
  \textbf{Supuestos:}
\end{enumerate}

\begin{itemize}
\item
  Se asume que el taller mantiene activa una conexión de red local
  (Wi-Fi) para la sincronización entre computadoras de mostrador y
  tabletas de taller en condiciones normales.
\item
  Se asume que los clientes cuentan con un correo electrónico para
  recibir el enlace de seguimiento y los códigos de verificación de 6
  dígitos (2FA); el número celular se usa para el envío opcional del
  enlace por SMS (RF-02.3).
\item
  Se asume que los datos ingresados durante la fase de desarrollo y
  pruebas son 100\% ficticios para dar cumplimiento a la Ley Federal de
  Protección de Datos Personales.
\end{itemize}

\begin{enumerate}
\def\labelenumi{\arabic{enumi}.}
\setcounter{enumi}{2}
\item
  \textbf{Dependencias:}
\end{enumerate}

\begin{itemize}
\item
  Infraestructura Eléctrica y Red: Operación dependiente de baterías en
  laptops/tabletas y arquitectura \emph{Local-First / PWA} (Service
  Workers e IndexedDB) para permitir la captura de datos \emph{offline}
  durante cortes de luz o fallas de internet.
\item
  Proveedores Externos de Comunicación: Dependencia del servicio de
  correo electrónico para entregar enlaces y códigos 2FA al cliente, y
  de un servicio de SMS para el envío opcional del enlace.
\item
  Disponibilidad de Refaccionarias Locales: La adquisición de piezas
  fuera de stock depende del directorio de proveedores locales e
  historial de compras del almacén.
\end{itemize}
